\section{Topologie \& Grundlagen im $\mathbb{R}^n$}

\subsection{Geometrische Grundlagen im $\mathbb{R}^n$}

\begin{definition}
    \textbf{Euklidischer Raum $\mathbb{R}^n$:} Der Raum $\mathbb{R}^n = \{(x_1, \dots, x_n)^T \mid x_i \in \mathbb{R}\}$ ist ein $n$-dimensionaler reeller Vektorraum mit der kanonischen Standardbasis $\{e_1, \dots, e_n\}$, wobei $e_i = (0, \dots, 1, \dots, 0)^T$ an der $i$-ten Stelle eine 1 besitzt.
\end{definition}

\begin{definition}
    \textbf{Standardskalarprodukt:} \\
    Für Vektoren $x, y \in \mathbb{R}^n$ ist das Standard-Skalarprodukt definiert durch:
    $$ \langle x, y \rangle := x^T y = \sum_{i=1}^n x_i y_i $$
    Es erfüllt die Eigenschaften:
    \begin{itemize}
        \item \textbf{Bilinearität:} $\langle \alpha x + \beta y, z \rangle = \alpha \langle x, z \rangle + \beta \langle y, z \rangle$ für alle $\alpha, \beta \in \mathbb{R}$.
        \item \textbf{Symmetrie:} $\langle x, y \rangle = \langle y, x \rangle$.
        \item \textbf{Positive Definitheit:} $\langle x, x \rangle \ge 0$ und $\langle x, x \rangle = 0 \iff x = 0$.
    \end{itemize}
\end{definition}

\begin{definition}
    \textbf{Euklidische Norm \& Euklidischer Abstand:}
    \begin{itemize}
        \item Die durch das Skalarprodukt induzierte \textbf{Euklidische Norm} (Länge) eines Vektors $x \in \mathbb{R}^n$ ist:
        $$ \|x\| := \sqrt{\langle x, x \rangle} = \sqrt{\sum_{i=1}^n x_i^2} $$
        \item Der \textbf{Euklidische Abstand} zweier Punkte $x, y \in \mathbb{R}^n$ ist die Metrik:
        $$ d(x, y) := \|x - y\| = \sqrt{\sum_{i=1}^n (x_i - y_i)^2} $$
    \end{itemize}
\end{definition}

\begin{theorem}
    \textbf{Cauchy-Schwarz-Ungleichung:} Für alle $x, y \in \mathbb{R}^n$ gilt:
    $$ |\langle x, y \rangle| \le \|x\| \cdot \|y\| $$
    Gleichheit gilt genau dann, wenn $x$ und $y$ linear abhängig sind (d.h. $\exists \lambda \in \mathbb{R}$ mit $x = \lambda y$ oder $y = 0$).
\end{theorem}

\begin{theorem}
    \textbf{Dreiecksungleichung:} Für alle $x, y, z \in \mathbb{R}^n$ gilt:
    $$ \|x + y\| \le \|x\| + \|y\| \quad \text{bzw.} \quad \|x - z\| \le \|x - y\| + \|y - z\| $$
    Zudem gilt die umgekehrte Dreiecksungleichung: $\big| \|x\| - \|y\| \big| \le \|x - y\|$.
\end{theorem}

\begin{definition}
    \textbf{$p$-Normen und Maximumsnorm:} Für $x \in \mathbb{R}^n$ definiert man:
    \begin{itemize}
        \item $p$-Norm ($1 \le p < \infty$): $\|x\|_p := \left( \sum_{i=1}^n |x_i|^p \right)^{1/p}$.
        \item Maximumsnorm ($p = \infty$): $\|x\|_\infty := \max_{1 \le i \le n} |x_i|$.
    \end{itemize}
\end{definition}

\begin{theorem}
    \textbf{Satz von der Normenäquivalenz im $\mathbb{R}^n$:} Auf dem endlichdimensionalen Vektorraum $\mathbb{R}^n$ sind alle Normen zueinander äquivalent. Das heisst: Zu zwei beliebigen Normen $\|\cdot\|_a$ und $\|\cdot\|_b$ auf $\mathbb{R}^n$ existieren Konstanten $c, C > 0$, sodass für alle $x \in \mathbb{R}^n$ gilt:
    $$ c \|x\|_a \le \|x\|_b \le C \|x\|_a $$
    \textbf{Konsequenz für Prüfungen:} Konvergenz von Folgen, Offenheit, Abgeschlossenheit und Kompaktheit von Mengen sind völlig unabhängig von der gewählten Norm im $\mathbb{R}^n$.
\end{theorem}

\subsection{Folgen im $\mathbb{R}^n$ und Konvergenz}

\begin{definition}
    \textbf{Konvergenz einer Vektorfolge:} Eine Folge $(x_k)_{k \in \mathbb{N}} \subset \mathbb{R}^n$ konvergiert gegen den Grenzwert $x^* \in \mathbb{R}^n$ (notiert als $\lim_{k \to \infty} x_k = x^*$ oder $x_k \to x^*$), falls:
    $$ \forall \varepsilon > 0 \; \exists N \in \mathbb{N} \; \forall k > N: \|x_k - x^*\| < \varepsilon $$
    Existiert kein solcher Grenzwert, heisst die Folge divergent.
\end{definition}

\begin{lemma}
    \textbf{Eigenschaften von Vektorfolgen:}
    \begin{enumerate}
        \item \textbf{Eindeutigkeit:} Der Grenzwert einer konvergenten Folge im $\mathbb{R}^n$ ist eindeutig.
        \item \textbf{Beschränktheit:} Jede konvergente Folge im $\mathbb{R}^n$ ist beschränkt (d.h. $\exists M > 0 \; \forall k: \|x_k\| \le M$).
        \item \textbf{Koordinatenweise Konvergenz:} Eine Vektorfolge $x_k = (x_{1,k}, \dots, x_{n,k})^T$ konvergiert genau dann im $\mathbb{R}^n$ gegen $x^* = (x_1^*, \dots, x_n^*)^T$, wenn jede Koordinatenfolge in $\mathbb{R}$ konvergiert:
        $$ \lim_{k \to \infty} x_k = x^* \iff \forall i \in \{1, \dots, n\}: \lim_{k \to \infty} x_{i,k} = x_i^* $$
    \end{enumerate}
\end{lemma}

\begin{theorem}
    \textbf{Satz von Bolzano-Weierstraß im $\mathbb{R}^n$:} Jede beschränkte Folge $(x_k)_{k \in \mathbb{N}}$ im $\mathbb{R}^n$ besitzt mindestens eine konvergente Teilfolge.
\end{theorem}

\begin{remark}
    \textbf{Beweisidee:} Wende den 1D-Satz von Bolzano-Weierstraß sukzessive $n$-mal an: Erst Teilfolge wählen, die in der 1. Koordinate konvergiert, daraus Teilfolge für die 2. Koordinate usw. Nach $n$ Schritten konvergieren alle $n$ Koordinaten.
\end{remark}

\subsection{Topologische Grundbegriffe im $\mathbb{R}^n$}

\begin{definition}
    \textbf{Offene und abgeschlossene Bälle:} Für $x_0 \in \mathbb{R}^n$ und Radius $r > 0$ definiert man:
    \begin{itemize}
        \item \textbf{Offener Ball:} $B_r(x_0) := \{x \in \mathbb{R}^n \mid \|x - x_0\| < r\}$.
        \item \textbf{Abgeschlossener Ball:} $\overline{B}_r(x_0) := \{x \in \mathbb{R}^n \mid \|x - x_0\| \le r\}$.
        \item \textbf{Sphäre (Kugeloberfläche):} $S_r(x_0) := \{x \in \mathbb{R}^n \mid \|x - x_0\| = r\}$.
    \end{itemize}
\end{definition}

\begin{definition}
    \textbf{Offene und abgeschlossene Mengen:} Sei $A \subset \mathbb{R}^n$.
    \begin{itemize}
        \item $A$ heisst \textbf{offen}, wenn jeder Punkt von $A$ ein innerer Punkt ist:
        $$ \forall x \in A \; \exists r > 0: B_r(x) \subset A $$
        \item $A$ heisst \textbf{abgeschlossen}, wenn ihr Komplement $A^c = \mathbb{R}^n \setminus A$ offen ist.
    \end{itemize}
\end{definition}

\begin{important}
    \textbf{Offen vs. Abgeschlossen:} Mengen sind keine Türen! Eine Menge kann offen, abgeschlossen, keines von beidem (z.B. $[0, 1) \subset \mathbb{R}$) oder beides gleichzeitig sein (in $\mathbb{R}^n$ sind nur $\emptyset$ und $\mathbb{R}^n$ sowohl offen als auch abgeschlossen, sogenannte \textit{clopen sets}).
\end{important}

\begin{lemma}
    \textbf{Schnitt- und Vereinigungseigenschaften:}
    \begin{itemize}
        \item Beliebige (auch überabzählbare) Vereinigungen offener Mengen sind offen.
        \item Endliche Durchschnitte offener Mengen sind offen.
        \item Beliebige Durchschnitte abgeschlossener Mengen sind abgeschlossen.
        \item Endliche Vereinigungen abgeschlossener Mengen sind abgeschlossen.
    \end{itemize}
\end{lemma}

\begin{definition}
    \textbf{Inneres, Abschluss und Rand:} Für eine Teilmenge $A \subset \mathbb{R}^n$ definiert man:
    \begin{itemize}
        \item \textbf{Inneres ($A^\circ$ bzw. $\interior(A)$):} Die Menge aller inneren Punkte, d.h. die grösste offene Teilmenge von $A$ (Vereinigung aller offenen Teilmengen von $A$).
        \item \textbf{Abschluss ($\overline{A}$ bzw. $\operatorname{cl}(A)$):} Der Schnitt aller abgeschlossenen Mengen, die $A$ enthalten (kleinste abgeschlossene Obermenge von $A$).
        \item \textbf{Topologischer Rand ($\partial A$):} Die Differenz $\partial A := \overline{A} \setminus A^\circ$.
    \end{itemize}
\end{definition}

\begin{lemma}
    \textbf{Charakterisierung über den Rand:}
    \begin{itemize}
        \item $A$ ist offen $\iff A \cap \partial A = \emptyset$ (kein Randpunkt gehört zu $A$).
        \item $A$ ist abgeschlossen $\iff \partial A \subset A$ (alle Randpunkte gehören zu $A$).
        \item Es gilt stets die disjunkte Zerlegung des Gesamtraums: $\mathbb{R}^n = A^\circ \;\dot{\cup}\; \partial A \;\dot{\cup}\; (\mathbb{R}^n \setminus \overline{A})$.
    \end{itemize}
\end{lemma}

\begin{theorem}
    \textbf{Folgencharakterisierung von Topologie:}
    \begin{itemize}
        \item \textbf{Abgeschlossenheit:} $A \subset \mathbb{R}^n$ ist abgeschlossen genau dann, wenn für jede Folge $(x_k)_{k \in \mathbb{N}} \subset A$, die im $\mathbb{R}^n$ gegen einen Grenzwert $x^*$ konvergiert, bereits $x^* \in A$ gilt (Abgeschlossenheit unter Grenzwertbildung).
        \item \textbf{Offenheit:} $A \subset \mathbb{R}^n$ ist offen genau dann, wenn für jedes $x \in A$ und jede Folge $x_k \to x$ ein Index $N \in \mathbb{N}$ existiert mit $x_k \in A$ für alle $k > N$.
    \end{itemize}
\end{theorem}

\subsection{Kompaktheit im $\mathbb{R}^n$}

\begin{definition}
    \textbf{Begriffe der Kompaktheit:}
    \begin{itemize}
        \item \textbf{Folgenkompaktheit:} Eine Teilmenge $K \subset \mathbb{R}^n$ heisst folgenkompakt, wenn jede Folge $(x_k)_{k \in \mathbb{N}} \subset K$ eine Teilfolge besitzt, die gegen einen Punkt in $K$ konvergiert.
        \item \textbf{Überdeckungskompaktheit:} $K \subset \mathbb{R}^n$ heisst kompakt, wenn jede offene Überdeckung $\{U_i\}_{i \in I}$ von $K$ (d.h. $K \subset \bigcup_{i \in I} U_i$ mit $U_i$ offen) eine endliche Teilüberdeckung besitzt: $\exists i_1, \dots, i_m \in I: K \subset U_{i_1} \cup \dots \cup U_{i_m}$.
    \end{itemize}
\end{definition}

\begin{remark}
    \textbf{Geometrische Anschauung: Warum Überdeckungskompaktheit? (Klausurklassiker)} \\
    \textbf{Warum scheitert das offene Intervall $(0, 1)$?} \\
    Die formale Forderung nach einer \textit{endlichen} Teilüberdeckung wirkt zunächst abstrakt. Ihre geometrische Kraft versteht man am besten an Mengen, die \textbf{nicht} kompakt sind.
    
    Betrachte das offene Intervall $K = (0, 1) \subset \mathbb{R}$ und die Familie offener Mengen
    $$ U_n := \left( \frac{1}{n}, 1 \right) \quad \text{für } n \ge 2 $$
    \begin{itemize}
        \item \textbf{Die Mengen bilden eine offene Überdeckung:} Für jedes noch so kleine $x \in (0, 1)$ existiert nach dem archimedischen Axiom ein $n \in \mathbb{N}$ mit $\frac{1}{n} < x$. Also gilt tatsächlich $(0, 1) \subset \bigcup_{n=2}^\infty U_n$.
        \item \textbf{Keine endliche Teilüberdeckung reicht aus:} Wählt man endlich viele Indizes $\{n_1, \dots, n_m\}$ und setzt $N := \max(n_1, \dots, n_m)$, so gilt:
        $$ U_{n_1} \cup \dots \cup U_{n_m} = \left( \frac{1}{N}, 1 \right) $$
        Jeder Punkt $x \in \left( 0, \frac{1}{N} \right]$ (wie etwa $x = \frac{1}{2N}$) wird von dieser endlichen Auswahl \textbf{nicht überdeckt}!
    \end{itemize}
    \textbf{Der geometrische Grund:} Weil der Randpunkt $0$ nicht zu $(0, 1)$ gehört, können sich die offenen Mengen immer näher an $0$ heranschleichen, ohne den Punkt jemals zu erfassen. Bei einem abgeschlossenen Intervall $[0, 1]$ müsste mindestens eine offene Menge den Punkt $0$ enthalten und würde damit automatisch eine ganze Umgebung $[0, \varepsilon)$ mit abdecken – dieses "Entwischen" wäre unmöglich!

    \textbf{Kernaussage für Prüfungen (Satz von Heine-Borel):} \\
    Im endlichdimensionalen Raum $\mathbb{R}^n$ fallen alle drei zentralen Kompaktheitsbegriffe zusammen:
    $$ \text{\textbf{Überdeckungskompakt}} \iff \text{\textbf{Folgenkompakt}} \iff \text{\textbf{Abgeschlossen \& Beschränkt}} $$
    In Klausuren beweist man Kompaktheit im $\mathbb{R}^n$ fast immer über \textbf{"abgeschlossen und beschränkt"}!
\end{remark}

\begin{theorem}
    \textbf{Satz von Heine-Borel:} Für eine Teilmenge $K \subset \mathbb{R}^n$ sind äquivalent:
    \begin{enumerate}
        \item $K$ ist kompakt (überdeckungskompakt).
        \item $K$ ist folgenkompakt.
        \item $K$ ist \textbf{abgeschlossen} und \textbf{beschränkt} (d.h. $\exists M > 0: \|x\| \le M$ für alle $x \in K$).
    \end{enumerate}
\end{theorem}